\begin{figure}[t]
    \label{fig:main}
    \centering
    \includegraphics[width=0.95\textwidth]{fig/main.pdf} 
    \caption{Comparison of story frames generation. The frame sequences are generated by various models, including (a)~StoryAnchors, (b)~StoryDiffusion, (c)~One Prompt One Story, (d)~VideoGen-of-Thought,  and (e)~GPT-4o.}
    \label{fig:main}
\end{figure}

\section{Experiments}
\label{sec:exp}
\subsection{Implement Details}
To evaluate the effectiveness of our proposed approach, we constructed a benchmark dataset consisting of 42 multi-event stories. The dataset was generated using Claude and further refined through manual annotation. 
It is designed to evaluate the model's ability to maintain consistency across multiple events and assess its capacity to generate coherent narratives and exhibit diversity throughout a sequence of story frames.

We adopt the Step-Video-T2V-4B variant~\cite{ma2025stepvideot2vtechnicalreportpractice,huang2025stepvideoti2vtechnicalreportstateoftheart} as our base model, with its text encoder built upon Step-LLM~\cite{ma2025stepvideot2vtechnicalreportpractice}. For the VAE component, we employ WanVAE~\cite{wan2025wanopenadvancedlargescale}.
During the annotation, keyframes were uniformly extracted from the original videos at intervals of 69 or 139 frames. An initial subset of the data was annotated using the Gemini system. 
Each video is center-cropped and resized to 480×768. The Adam optimizer with a learning rate of 1e-5 is used for each training stage. A total batch size of 256 is distributed across 64 H100 GPUs.


\subsection{Comparisons of Consistent Story Frames Generation}

\textbf{Qualitative Comparisons.} 
Figure~\ref{fig:main} presents a qualitative comparison between our method (\ours{}) and several baselines for consistent story frame generation, including StoryDiffusion~\cite{zhou2024storydiffusion}, VideoGen-of-Thought~\cite{zheng2025videogenofthoughtstepbystepgeneratingmultishot}, One Prompt One Story~\cite{liu2025one}, and GPT-4o. Each sequence illustrates how a model interprets a given prompt to construct a coherent visual narrative.
Among the methods evaluated, our model, along with VideoGen-of-Thought and One Prompt One Story, consistently produces sequences with higher thematic and background coherence. Notably, \ours{} demonstrates greater shot diversity, featuring a broader range of camera angles that enhance visual storytelling. In contrast, GPT-4o often defaults to mid-shots, resulting in limited compositional variety and reduced narrative engagement.
While StoryDiffusion generates relatively diverse shots, it struggles with maintaining logical event progression and temporal consistency, particularly in complex scenes. In comparison, our model exhibits stronger narrative coherence and a clearer depiction of event dynamics and emotional transitions.
Moreover, GPT-4o frequently suffers from detail inconsistencies and color shifts, especially as the number of generated keyframes or dialogue turns increases. Our model, by contrast, maintains both visual fidelity and consistency across frames, while delivering comparable or superior performance in narrative quality and scene richness.

\begin{figure}[t]
    \centering
    \subfigure[One Prompt One Story~(10)]{\includegraphics[width=0.3\linewidth]{result_fig/OnePromptOneStory.pdf}}
    \quad
    \subfigure[StoryDiffusion~(10)]{\includegraphics[width=0.3\linewidth]{result_fig/StoryDiffusion.pdf}}  
    \quad
    \subfigure[VideoGen of Thought~(10)]{\includegraphics[width=0.3\linewidth]{result_fig/VideoGenofThought.pdf}}  
    \\
   \subfigure[One Prompt One Story~(20)]{\includegraphics[width=0.3\linewidth]{result_fig/OnePromptOneStory20.pdf}}
    \quad
    \subfigure[StoryDiffusion~(20)]{\includegraphics[width=0.3\linewidth]{result_fig/StoryDiffusion20.pdf}} 
    \quad
    \subfigure[VideoGen of Thought~(20)]{\includegraphics[width=0.3\linewidth]{result_fig/VideoGenofThought20.pdf}}  
    \caption{Comparison of Open-Source Methods Across Five Evaluation Metrics (Evaluated by Gemini).
This table summarizes the performance of \ours{} and other baseline methods under five evaluation metrics assessed by the Gemini system. The values in parentheses indicate the number of frames generated by each method.}
    \label{fig:2x3_subfigures}
\end{figure}
\vspace{8pt}

\textbf{Quantitative Comparisons.} For the quantitative evaluation, we adopt five keyframe-based automatic metrics, inspired by  prior works~\cite{ zhou2024storydiffusion,zheng2025videogenofthoughtstepbystepgeneratingmultishot}: Narrative Quality~(NQ), Subject Consistency~(SC), Prompt-Frame Alignment~(PFA), Shot Diversity~(ShotD), and Scene Diversity~(SceneD). A detailed explanation of each metric is provided in Appendix~\ref{app:metrics}.

\begin{wraptable}{r}{0.55\textwidth}
\vspace{-3.5mm}
\centering
\small 
\caption{Comparison results (Win/Tie/Loss) based on story frames generated by GPT-4o.}
\begin{tabular}{@{}lcccccc@{}}
\toprule
\textbf{Metric} & \textbf{NQ} & \textbf{SC} & \textbf{PFA} & \textbf{ShotD} & \textbf{SceneD} & \textbf{Total} \\
\midrule
Win  & 7  & 10 & 5  & 7  & 10 & 39 \\
Tie  & 20 & 18 & 24 & 29 & 22 & 113 \\
Loss & 15 & 14 & 13 & 6  & 10 & 58 \\
\bottomrule
\end{tabular}
\end{wraptable}
The quantitative results are in line with our qualitative findings. When generating 10 frames, \ours{} achieves strong overall performance, particularly excelling in subject consistency and scene diversity. For other open-source methods, increasing the number of generated keyframes tends to enhance narrative richness and variation. However, this improvement often comes at the cost of reduced consistency, as these models struggle to maintain coherent subject representations and stable scene layouts beyond the 10-frame mark.
We also conduct a comparative evaluation against GPT-4o. As shown in Table~\ref{tab:gpt4o_single_turn_transposed}, GPT-4o meets the resolution and instruction-following requirements in 71.4\% of cases. While \ours{} yields slightly lower scores in NQ and PFA, it delivers comparable overall results, with several metrics showing parity between the two models.

\begin{figure}[t]
    \centering
    \includegraphics[width=0.9\textwidth]{fig/edit1.pdf} 
    \caption{The generated frames are conditioned on the first four frames as the additional condition, and different event prompts are used to guide the continuation. }
    \label{fig:ed}
\end{figure}

\begin{figure}[t]
\centering
\includegraphics[width=\linewidth]{fig/edit2.pdf} % 根据实际路径修改
\caption{
The model is conditioned on different frames (2 Frames, 3 Frames, and 4 Frames) and generates the following sequence based on the same multi-event caption. 
}
\label{fig:frame_conditioning}
\end{figure}

\subsection{Story Frames Editing and Temporal Extension}

In Figure~\ref{fig:ed}, we further validate this property by conditioning the model on the first five frames of a video while using different prompts to guide continuation. The results demonstrate that the model successfully preserves key subject features observed within the conditioned frames while generating diverse and story-consistent frames based on the varying prompts.
Interestingly, although the model was trained with only the first frame as input, it exhibits strong generalization capabilities at inference time. We observe that conditioning the model on a single frame is sufficient for learning meaningful temporal dynamics. Moreover, when provided with multiple frames during inference, the model can still effectively perform event continuation and maintain narrative coherence, despite never being explicitly trained under multi-frame conditions.

Additionally, we conduct an experiment where the same model is conditioned on different single frames and then tasked with generation under the same prompt. As shown in Figure~\ref{fig:frame_conditioning}, across different conditioning frames, the model consistently produces plausible continuations, maintaining both subject identity and narrative consistency. These findings highlight the model’s flexibility and its robustness in handling variations in the conditioning inputs,



\subsection{Effectiveness of Multi-Event Prompt}
\begin{figure}[t]
    \centering
        \includegraphics[width=0.95\linewidth]{fig/ablation.crop.pdf}
    \caption{Comparison between global prompt and multi-event descriptions.
Each subfigure presents a generated video sequence using (top) a single global prompt and (bottom) multi-event prompt. }
    \label{fig:full}
\end{figure}

The \textbf{Global Prompt} method treats the video as a single, continuous sequence and generates frames based on a high-level description without distinguishing individual events. This approach results in significant challenges in maintaining both \textbf{event completeness} and \textbf{temporal order}. For instance, in Figure~\ref{fig:prompt_comparison}(a), the global-prompt-based generation omits important moments such as the squirrel jumping over the stone and the girl’s walk along the lakeside, both of which are crucial for a complete depiction of the narrative. Similarly, in Figure~\ref{fig:prompt_comparison}(b), the global prompt introduces large numbers of butterflies too early and misplaces the butterfly-shaped map at the end, despite it being a key discovery that initiates the journey, thereby disrupting the logical event progression.

In contrast, our \textbf{Multi-Event Description} approach explicitly segments the story into structured events, each guided by detailed descriptions. This enables more accurate and coherent visual storytelling, ensuring \textbf{narrative consistency}, \textbf{event completeness}, and \textbf{temporally aligned} generation. As shown in both examples, the multi-event method captures all significant moments in the correct sequence, resulting in a more faithful and interpretable depiction of the story.

